\chapter{DISCUSIÓN}

Este capítulo integra los resultados de los dos ejes experimentales y examina sus implicaciones para el diseño de sistemas AML auditables. La discusión parte de la lectura conjunta de Elliptic y del grafo sintético, deriva de ella una reconsideración de qué factores gobiernan la calidad de las explicaciones, y sitúa los hallazgos frente al estado del arte. El capítulo cierra con una exposición honesta de las limitaciones del estudio, dado que reconocer sus fronteras es condición para interpretar correctamente su alcance.

\section{Lectura Conjunta de los Dos Ejes}

La combinación de los dos ejes ofrece una imagen más completa que la que cualquiera de ellos podría dar por separado, y esa complementariedad es precisamente el argumento metodológico central de la tesis. El eje Elliptic aporta el realismo de datos de Bitcoin con ruido, desbalance y desplazamiento temporal, y en él se establece, una vez corregidos dos artefactos de evaluación sucesivos y replicado el entrenamiento con tres semillas de modelo, que la estabilidad no ordena a las cuatro arquitecturas en cuatro posiciones distinguibles sino que las separa en dos grupos, uno alto formado por GAT y GCN y uno bajo formado por GraphSAGE y TAGCN, con diferencias significativas entre grupos y no dentro de ellos. El eje sintético aporta el control necesario para medir la plausibilidad, y en él se establece que el explicador es el factor determinante, que PGExplainer recupera mejor el patrón de referencia y que no existe un puente entre estabilidad y plausibilidad. Ninguno de los dos ejes contradice al otro, sino que cada uno responde preguntas que el otro no puede abordar, de modo que su lectura conjunta sostiene conclusiones más robustas que las de un diseño de un solo dataset.

Un punto de contacto entre ambos ejes merece subrayarse, el comportamiento de PGExplainer. En Elliptic, PGExplainer degeneraba de forma casi universal y producía correlaciones nulas, lo que en un estudio de un solo dataset habría llevado a concluir que se trata de un método deficiente. El eje sintético desmiente esa conclusión, ya que sobre un grafo denso PGExplainer se convierte en el explicador más plausible, lo que revela que su fracaso en Elliptic era una consecuencia de la dispersión del dato y no una propiedad del método. Este contraste ilustra el valor de disponer de dos regímenes de densidad distintos, puesto que permite distinguir entre las limitaciones de un método y las limitaciones que el dato impone a cualquier método.

\section{Estabilidad, Plausibilidad y Fidelidad como Tres Dimensiones Independientes}

El hallazgo más consecuente de la tesis es que la estabilidad, la plausibilidad y la fidelidad de una explicación son dimensiones independientes que deben medirse por separado. El puente nulo entre estabilidad y plausibilidad demuestra que un explicador consistente entre ejecuciones no acierta por ello sobre el patrón real, y la disociación entre plausibilidad y fidelidad demuestra que un explicador que recupera el patrón que un humano reconoce no refleja necesariamente el mecanismo interno del modelo. Estas dos relaciones, tomadas en conjunto, desmontan la intuición de que existe una única noción de buena explicación que las tres métricas capturarían de forma redundante. Por el contrario, cada dimensión responde a una pregunta distinta, la estabilidad pregunta si la explicación es reproducible, la plausibilidad pregunta si coincide con el conocimiento del dominio, y la fidelidad pregunta si describe lo que el modelo hace.

La independencia de estas tres dimensiones sugiere que la evaluación de explicaciones debería adoptar un enfoque multidimensional de forma sistemática. Reportar una única métrica de calidad, como es habitual en muchos trabajos, oculta compromisos que pueden ser decisivos para la aplicación. Un método que encabeza un ranking de estabilidad podría ocupar el último lugar en fidelidad, de modo que una comparación basada en una sola dimensión conduciría a una recomendación equivocada según el propósito real de la explicación. Esta tesis aboga, en consecuencia, por caracterizar cada explicador a través de las tres dimensiones de forma conjunta, y por hacer explícito el propósito de la auditoría antes de seleccionar un método.

Esta independencia tiene consecuencias directas para la práctica de la auditoría en entornos regulados. Un equipo de cumplimiento que necesite justificar una alerta ante un regulador debe decidir qué dimensión prioriza, ya que ningún explicador las optimiza todas a la vez. Si el objetivo es demostrar que la alerta responde a un patrón de lavado reconocible, la plausibilidad es la dimensión relevante y PGExplainer resulta preferible, aunque su baja fidelidad advierta que la explicación puede no coincidir con el razonamiento del modelo. Si el objetivo es auditar el modelo en sí, para verificar en qué se basa realmente su decisión, la fidelidad es la dimensión relevante y GNNExplainer resulta preferible, aunque su plausibilidad sea menor. La elección, por tanto, no es técnica sino que depende del propósito de la auditoría, una conclusión que la literatura sobre explicabilidad rara vez plantea de forma tan explícita \cite{Agarwal2022ProbingMethods, Adadi2018PeekingXAI}.

\section{El Papel Secundario del Balanceo y de la Arquitectura}

Un resultado que contrasta con parte de la literatura es el papel prácticamente nulo del balanceo sobre la calidad de las explicaciones. Trabajos previos habían sugerido que las estrategias de mitigación del desbalance influyen tanto en el clasificador como en la consistencia de las explicaciones \cite{Rai2025EvaluatingData, Fujiwara2024KnowledgeStability}. Los resultados del eje sintético, en cambio, muestran que el balanceo tiene un tamaño de efecto despreciable sobre la estabilidad y la plausibilidad, muy por debajo del efecto del explicador. Esta discrepancia puede explicarse porque los trabajos anteriores median el efecto del balanceo sobre el clasificador o sobre datos distintos de los grafos financieros, mientras que la presente tesis lo mide directamente sobre la calidad de las explicaciones en un grafo con \textit{ground-truth}, un encuadre que aísla mejor la contribución de cada factor.

La arquitectura ocupa una posición intermedia, con un efecto apreciable sobre la estabilidad pero pequeño sobre la plausibilidad. Una vez corregidos los dos artefactos de evaluación, el orden de estabilidad entre arquitecturas deja de invertirse entre Elliptic y el grafo sintético y pasa a ser concordante, con GAT y GCN a la cabeza en ambos regímenes. La replicación con tres semillas precisa el alcance de esta concordancia y la refuerza, ya que los dos regímenes no solo coinciden en qué arquitecturas encabezan, sino en la partición completa, con GAT y GCN separados de GraphSAGE y TAGCN. La separación entre los dos grupos alcanza significación estadística en ambos ejes, y dentro del grupo bajo las dos arquitecturas resultan indistinguibles también en ambos, con valores p de 0,096 en Elliptic y de 0,519 en el grafo sintético. Dentro del grupo alto la distinción depende del régimen, puesto que en Elliptic GAT y GCN no se separan mientras que en el grafo sintético GCN supera a GAT con un valor p de 0,0013, si bien por un margen de seis milésimas que carece de relevancia práctica y que solo resulta detectable por el mayor número de celdas y la menor varianza de ese eje. Que dos regímenes de densidad opuesta, medidos sobre datos distintos y con pipelines independientes, converjan en la misma partición ofrece un respaldo más fuerte que la coincidencia de un ordenamiento entre cuatro posiciones. La conclusión práctica se modifica en consecuencia, puesto que la elección de arquitectura sí ofrece una palanca de estabilidad, aunque de grano grueso, ya que lo que la evidencia sostiene es la elección de grupo y no la elección de una arquitectura concreta dentro de él.

La concordancia del orden de estabilidad entre regímenes tiene una implicación teórica que conviene explicitar, y la partición en dos grupos la vuelve más nítida de lo que permitía un ordenamiento de cuatro posiciones. La composición de los grupos no es arbitraria desde el punto de vista del mecanismo de agregación. GCN y GAT, las dos arquitecturas del grupo alto, ponderan la vecindad inmediata de forma determinista una vez fijados los parámetros, la primera mediante una normalización espectral fija y la segunda mediante pesos de atención que, aunque aprendidos, no varían en inferencia. GraphSAGE y TAGCN, las dos del grupo bajo, introducen en cambio sendas fuentes de variabilidad, ya que la primera muestrea la vecindad y la segunda reparte la importancia sobre un campo receptivo ampliado por el orden del filtro polinomial. La estabilidad de una explicación dependería, según esta lectura, de si el mecanismo de agregación concentra la atribución sobre un conjunto reducido y fijo de vecinos o la dispersa sobre un conjunto amplio o variable. Aunque esta interpretación requiere confirmación adicional, el hecho de que la misma partición emerja de forma independiente en dos regímenes de densidad opuesta le otorga un respaldo apreciable, y abre una línea de investigación sobre la relación entre la agregación y la estabilidad explicativa.

\section{Contribuciones Metodológicas}

Más allá de los hallazgos sustantivos, la tesis aporta varias contribuciones metodológicas que conviene destacar. La primera es la corrección del subgrafo receptivo, que revela cómo un fallo silencioso de memoria puede confundirse con una propiedad de una arquitectura, un riesgo general en la evaluación de explicadores que rara vez se documenta. La segunda es la distinción entre la degeneración de un método y la limitación que el dato impone, ilustrada por el comportamiento de PGExplainer en los dos regímenes de densidad. La tercera es la medición uniforme de la fidelidad para los tres explicadores, que hace posible la comparación completa que dio lugar al hallazgo de la disociación, y que suple una carencia de las implementaciones estándar, las cuales no cubren por igual a todos los métodos. La cuarta es la distinción, establecida de forma empírica al replicar el entrenamiento, entre la reproducibilidad de los pesos y la reproducibilidad de las conclusiones. Reentrenar la matriz con la misma semilla no devuelve modelos idénticos, porque las operaciones de dispersión del paso de mensajes sobre la unidad de procesamiento gráfico emplean sumas atómicas de orden no determinista, y sin embargo devuelve el mismo filtro de calidad y el mismo orden de estabilidad. Enunciar esa distinción evita tanto una expectativa de reproducibilidad exacta que el cómputo sobre aceleradores gráficos no puede satisfacer, como la conclusión opuesta y errónea de que un resultado no reproducible bit a bit carece de valor científico.

A estas contribuciones se suma una práctica de honestidad metodológica que atraviesa todo el trabajo, la de retractar conclusiones anteriores cuando la evidencia corregida las contradice. La tesis retira de forma explícita el supuesto compromiso entre exactitud y estabilidad, el supuesto pico de inestabilidad en un escenario intermedio y la supuesta paradoja del escenario nativo, porque todos ellos se apoyaban en el artefacto de memoria ya descrito. Retractar hallazgos no es un signo de debilidad, sino una manifestación del rigor que exige sustituir conclusiones apoyadas en artefactos por conclusiones apoyadas en mediciones homogéneas. Esta actitud es coherente con las advertencias de la literatura sobre la sensibilidad de los resultados de explicabilidad al protocolo experimental \cite{Kosan2023GNNX-BENCH:Benchmarking}.

\section{Relación con el Estado del Arte}

Los hallazgos de esta tesis dialogan de forma directa con varios trabajos del estado del arte, a los que confirman en algunos aspectos y matizan en otros. La distinción entre plausibilidad y fidelidad que aquí se cuantifica coincide con la advertencia teórica de Agarwal y colaboradores, quienes señalaron que muchas explicaciones consideradas plausibles no son necesariamente fieles al mecanismo del modelo \cite{Agarwal2022ProbingMethods}. La contribución de esta tesis consiste en llevar esa advertencia del plano conceptual al plano de la medición, mostrando con un grafo de \textit{ground-truth} conocido que la plausibilidad y la fidelidad no solo difieren, sino que pueden ordenarse de manera opuesta entre explicadores. De este modo, la disociación deja de ser una posibilidad teórica y se convierte en un fenómeno observado y cuantificado sobre un conjunto de datos controlado.

La sensibilidad de los resultados al protocolo de evaluación, documentada por Kosan y colaboradores mediante un \textit{benchmark} exhaustivo, encuentra en esta tesis una ilustración concreta y de consecuencias graves \cite{Kosan2023GNNX-BENCH:Benchmarking}. Los dos artefactos que alteraron el ranking de estabilidad, primero un fallo de memoria que favorecía a GAT y después un truncamiento de la métrica de Spearman que favorecía a GraphSAGE, son precisamente el tipo de dependencia del protocolo que aquel trabajo advertía, y su corrección sucesiva demuestra que un detalle de implementación puede determinar por completo la conclusión de un estudio comparativo. Este episodio refuerza la necesidad de auditar tanto los métodos como las condiciones de cómputo bajo las que se miden. La coincidencia entre una advertencia general de la literatura y un caso concreto en este estudio otorga a la corrección un valor que trasciende el resultado puntual.

La relación con los trabajos sobre desbalance y explicabilidad merece un comentario cuidadoso. Rai y colaboradores junto con Fujiwara reportaron que el desbalance de clases afecta a la consistencia de las explicaciones \cite{Rai2025EvaluatingData, Fujiwara2024KnowledgeStability}. Los resultados de esta tesis no contradicen esa observación en cuanto al clasificador, pero sí matizan su alcance sobre la calidad de las explicaciones, ya que el balanceo, medido de forma aislada sobre un grafo con \textit{ground-truth}, ejerce un efecto despreciable frente al del explicador. Esta diferencia sugiere que el efecto del desbalance sobre las explicaciones podría estar mediado por su efecto sobre el clasificador, y no ser un efecto directo sobre el explicador, una hipótesis que trabajos futuros podrían examinar. Finalmente, la fragilidad de los detectores basados en grafos ante el desplazamiento temporal, señalada por Lawal y por Menezes junto con sus colaboradores, se manifiesta con nitidez en el colapso de validación a test documentado en el Capítulo 4 \cite{Lawal2025AnDataset, Menezes2025SemanticDetectors}.

\section{Implicaciones para la Práctica de Auditoría en Entornos Regulados}

Las conclusiones de esta tesis tienen implicaciones concretas para el diseño de sistemas de detección AML que deban ser auditables ante un regulador. La primera implicación es que la elección de un explicador no puede tratarse como una decisión puramente técnica, sino que debe alinearse con el propósito de la auditoría. Un equipo que necesite demostrar que una alerta responde a un patrón de lavado reconocible priorizará la plausibilidad, mientras que un equipo que deba certificar en qué se basa realmente el modelo priorizará la fidelidad, y estos dos objetivos conducen a explicadores distintos. Ignorar esta distinción puede llevar a presentar como justificación del modelo una explicación que en realidad no refleja su razonamiento, con el riesgo jurídico que ello implica.

La segunda implicación concierne a la estabilidad como requisito operativo. Una explicación que cambia de forma sustancial entre ejecuciones carece de valor probatorio, puesto que un analista no podría sostener ante un auditor una evidencia que no es reproducible. Los resultados indican que la estabilidad depende con fuerza del explicador y de la arquitectura, pero apenas del balanceo, lo que orienta las decisiones de diseño hacia la selección del par formado por la arquitectura y el explicador, y aleja la atención de la estrategia de balanceo como palanca de estabilidad. En el régimen disperso de los datos reales de Bitcoin, la combinación de una arquitectura de agregación determinista como GAT o GCN con un explicador estable como GNNShap ofrece el mejor punto de partida para obtener explicaciones reproducibles.

La tercera implicación es de carácter metodológico y afecta a la propia confianza en los sistemas de evaluación. El hecho de que sendos fallos silenciosos, uno de cómputo y otro de medida, pudieran invertir una conclusión sobre qué arquitectura es más estable advierte que las evaluaciones de explicabilidad deben incluir verificaciones de integridad tanto del cómputo como de las métricas, y no solo de los métodos. En un contexto regulado, donde las decisiones se apoyan en estas evaluaciones, la ausencia de tales verificaciones constituye un riesgo que esta tesis contribuye a visibilizar. La honestidad en la retractación de hallazgos apoyados en artefactos es, en este sentido, una práctica que los estudios de este tipo deberían adoptar de forma sistemática para preservar la confianza en sus conclusiones.

\section{Implicaciones Más Amplias del Estudio}

Los hallazgos de esta tesis trascienden el problema técnico concreto y tocan cuestiones de confianza, ética y regulación que conviene discutir de forma explícita. Las tres subsecciones siguientes exploran estas implicaciones más amplias, y muestran que la caracterización multidimensional de la calidad explicativa tiene consecuencias que van más allá de la elección de un método.

\subsection{Explicabilidad y Confianza en los Sistemas de Inteligencia Artificial}

La confianza que un analista o un regulador deposita en un sistema de detección no se sustenta únicamente en la exactitud de sus predicciones, sino también en la calidad de las razones que ofrece para justificarlas. Cuando un modelo alcanza un desempeño predictivo elevado pero acompaña cada alerta con explicaciones que cambian de una ejecución a otra, el usuario percibe que el sistema carece de un criterio consistente y termina desconfiando incluso de sus aciertos. Esta observación revela una distinción que resulta central para esta tesis: confiar en la predicción y confiar en la explicación son actos separados, gobernados por evidencias distintas. Un analista puede aceptar que una transacción es ilícita y, al mismo tiempo, rechazar la narrativa que el explicador construye alrededor de esa decisión, porque esa narrativa no le resulta ni estable ni verificable.

Los hallazgos de esta investigación refuerzan esa separación al mostrar que la estabilidad, la plausibilidad y la fidelidad constituyen dimensiones independientes de la calidad explicativa. Una explicación inestable erosiona la confianza porque impide que el analista construya un modelo mental duradero del comportamiento del sistema, ya que las mismas entradas producen justificaciones que fluctúan sin un patrón discernible. Una explicación infiel resulta aún más peligrosa, dado que ofrece al usuario una historia coherente y persuasiva que, sin embargo, no corresponde al razonamiento real del modelo. En ambos casos la precisión del clasificador queda subordinada a la credibilidad de su interfaz explicativa, porque el usuario humano actúa sobre lo que comprende y no sobre lo que el modelo calcula internamente.

Reconocer que estas tres propiedades no covarían tiene consecuencias directas sobre cómo se debe cultivar la confianza en entornos de auditoría financiera. Un explicador que produce mapas visualmente convincentes puede generar una confianza injustificada si esos mapas no reflejan las variables que efectivamente determinan la salida del modelo. Por el contrario, un explicador estable y fiel, aunque ofrezca explicaciones menos intuitivas, sienta una base más sólida para que el analista calibre cuándo delegar en el sistema y cuándo revisar manualmente. De este modo, la calidad de las explicaciones se convierte en un requisito para que la confianza sea legítima y no meramente aparente, y esta tesis argumenta que dicha calidad debe evaluarse de forma multidimensional en lugar de reducirse a una sola cifra agregada.

\subsection{Consideraciones Éticas de la Detección Automatizada}

Un sistema capaz de bloquear cuentas o de generar reportes con consecuencias legales opera sobre la vida de personas concretas, lo que traslada la discusión desde el plano técnico hacia el terreno de la responsabilidad ética. Cada alerta que el modelo produce puede desencadenar la congelación de fondos, la apertura de una investigación o la exclusión de un usuario del sistema financiero, medidas que afectan de manera desproporcionada a quienes resultan señalados por error. El costo humano de los falsos positivos no se distribuye de forma uniforme, puesto que las poblaciones con menor capacidad de reclamación suelen ser las que más difícilmente revierten una decisión automatizada. Por ello, el diseño de estos sistemas no puede tratar los falsos positivos como un simple parámetro a optimizar, sino como un daño real que exige mecanismos de reparación y revisión.

En este contexto, el derecho de una persona a recibir una explicación comprensible sobre por qué fue señalada adquiere un peso ético que trasciende la conveniencia operativa. Una explicación que el afectado y el auditor puedan interpretar es condición para que exista un debido proceso, dado que nadie puede impugnar una decisión cuyo fundamento permanece oculto o resulta ininteligible. La fidelidad de la explicación se vuelve entonces un asunto de justicia y no solo de calidad técnica, porque una justificación que no corresponde al razonamiento real del modelo priva al afectado de la posibilidad de defenderse frente a los motivos verdaderos de la decisión. Cuando la explicación es infiel, el proceso de apelación se construye sobre una ficción, y la garantía de defensa se transforma en una formalidad vacía.

A esta preocupación se suma el riesgo de que el modelo apoye sus decisiones en características espurias que correlacionan con la etiqueta sin guardar relación causal con la conducta ilícita. Si un clasificador aprende a asociar la ilicitud con atributos accidentales del conjunto de datos, puede penalizar sistemáticamente a usuarios que comparten esos atributos sin haber cometido ninguna irregularidad. La fidelidad explicativa cumple aquí una función de vigilancia, ya que solo una explicación que revele el razonamiento efectivo del modelo permite detectar cuándo este se apoya en señales indebidas. Por consiguiente, esta tesis sostiene que la evaluación de la fidelidad no es un ejercicio académico, sino una salvaguarda necesaria para que la automatización no reproduzca ni amplifique sesgos con impacto sobre personas inocentes.

\subsection{El Papel de la Explicabilidad en la Adopción Regulatoria de la IA}

La incorporación de modelos complejos en el sector financiero depende, en última instancia, de que los reguladores acepten decisiones cuyo mecanismo interno no siempre es transparente. Las autoridades de supervisión exigen que las entidades puedan justificar sus determinaciones ante auditorías, lo que coloca a la explicabilidad en el centro del proceso de adopción tecnológica. Existe, sin embargo, una tensión persistente entre el desempeño de los modelos más opacos y la exigencia regulatoria de transparencia, dado que las arquitecturas que capturan relaciones complejas suelen ser también las más difíciles de interpretar. Esta tensión explica por qué muchas instituciones prefieren modelos menos potentes pero legibles, sacrificando capacidad predictiva para preservar la posibilidad de rendir cuentas.

La contribución de esta tesis a ese debate consiste en mostrar que la calidad explicativa puede caracterizarse de forma estructurada a través de sus tres dimensiones, en lugar de tratarse como una propiedad monolítica que un método posee o no posee. Al distinguir la estabilidad, la plausibilidad y la fidelidad, el marco propuesto ofrece a los reguladores un vocabulario más preciso para especificar qué tipo de garantía explicativa requiere cada uso, ya que una auditoría orientada a la reparación de daños no demanda lo mismo que una orientada a la validación científica del modelo. Este enfoque también clarifica que ningún explicador optimiza simultáneamente las tres dimensiones, de modo que la selección del método debe subordinarse al propósito concreto de la supervisión. De esta manera, la elección técnica deja de ser arbitraria y pasa a responder a criterios que un regulador puede formular y verificar.

Caracterizar la calidad explicativa en estos términos ayuda a construir un puente entre la capacidad técnica de los modelos y su aceptación normativa, porque traduce una preocupación abstracta sobre la transparencia en requisitos evaluables y comparables. Cuando el regulador dispone de dimensiones separadas para exigir estabilidad, plausibilidad o fidelidad según el escenario, la conversación entre desarrolladores y supervisores puede apoyarse en criterios compartidos en lugar de en apreciaciones subjetivas. Esta claridad reduce la incertidumbre que frena la adopción, dado que una entidad sabe de antemano qué propiedad explicativa deberá demostrar para que un modelo complejo resulte admisible. En síntesis, esta tesis defiende que el avance hacia sistemas más potentes en el ámbito financiero no pasa por renunciar a la transparencia, sino por descomponerla en dimensiones que la técnica pueda entregar y que la regulación pueda auditar.


\section{Limitaciones}

El estudio presenta limitaciones que delimitan con honestidad su alcance y que deben tenerse presentes al interpretar sus conclusiones. La primera es que el eje real se apoya en un único conjunto de datos anonimizado, el \textit{Elliptic Dataset}, cuyo desplazamiento temporal provoca el colapso de validación a test y cuya anonimización impide medir la plausibilidad, de modo que las conclusiones sobre datos reales se restringen a la estabilidad sobre validación. La segunda es que, aun con la replicación por tres semillas de modelo en ambos ejes, cada celda de la matriz descansa en un número limitado de entrenamientos, y aunque la dispersión observada es baja en la mayoría de las arquitecturas, un número mayor de semillas estrecharía los intervalos de confianza. Esta limitación resulta especialmente pertinente en el caso de TAGCN sobre Elliptic, cuya dispersión entre semillas triplica la de las demás arquitecturas, de modo que su posición dentro del grupo bajo está bien establecida pero su valor central admite un margen considerable. La tercera es que el grafo sintético, por construcción, carece del desplazamiento temporal y del ruido de las transacciones reales, y sus tipologías son inyecciones controladas que no reproducen toda la variedad del lavado real, de modo que su validez es interna y no externa.

Una limitación adicional concierne a la comparación de fidelidad entre familias de explicadores. La fidelidad de GNNExplainer y PGExplainer se mide enmascarando aristas, mientras que la de GNNShap se mide enmascarando features, ya que cada método opera sobre un tipo distinto de máscara. Esta diferencia se declara de forma explícita y obliga a leer con cautela la comparación de fidelidad entre familias, aunque no afecta a la comparación dentro de una misma familia ni al hallazgo de la disociación, que se establece entre PGExplainer y GNNExplainer sobre aristas. Una consideración adicional concierne a la generalización de los hallazgos más allá del dominio AML. Aunque el estudio se sitúa en la detección de transacciones ilícitas, los tres fenómenos centrales, la independencia entre estabilidad, plausibilidad y fidelidad, la concordancia del orden de estabilidad entre regímenes de densidad del grafo, y la distinción entre la degeneración de un método y la limitación del dato, no son propios de las finanzas, sino que atañen a la explicabilidad de las redes neuronales sobre grafos en general. Es razonable esperar que se manifiesten en otros dominios que compartan las condiciones estructurales relevantes, como la dispersión de la red o la existencia de un patrón de referencia conocido. No obstante, esta generalización es una conjetura que requiere verificación en cada dominio, y no una conclusión que el presente estudio pueda afirmar por sí solo, dado que sus experimentos se circunscriben a los dos ejes descritos.

En conjunto, estas limitaciones no invalidan los hallazgos, pero sí definen el marco dentro del cual deben interpretarse, y varias de ellas señalan direcciones concretas de trabajo futuro.
